\Article{297}
Anyone looking more closely at the preceding proof will readily see
that cases I and II have been settled without anything to object to.
But the proofs of the other cases rely on the existence of auxiliary
numbers. Unless that existence is proved, the method clearly loses
all force. These assumptions may be so plausible that a less attentive
reader might think they need no proof, and they certainly raise the
theorem to be proved to the highest degree of \emph{probability};
still, if geometric rigor is wanted, they must not be admitted without
proof. As for the assumption in IV and V that there is a prime \(r\)
of the form \(4n+1\) that is a nonresidue of two other given primes
\(p,q\), Section IV readily shows that all numbers less than \(4pq\)
and coprime to it (there are \(2(p-1)(q-1)\) of them) divide equally into four
classes. One contains nonresidues of both \(p,q\); the other three,
residues of \(p\) and nonresidues of \(q\), nonresidues of \(p\)
and residues of \(q\), and residues of both \(p,q\). In each class,
half the numbers are of the form \(4n+1\) and half of the form \(4n+3\).
So among them there will be \(\frac14(p-1)(q-1)\) nonresidues of both \(p,q\)
of the form \(4n+1\). Call these \(g.g',g''\), and so on, and let the
remaining \(\frac74(p-1)(q-1)\) numbers be \(h,h',h''\), and so on. Clearly all numbers
contained in the forms \(4pqt+g,4pqt+g',4pqt+g''\), and so on, \((G)\) will also be
nonresidues of \(p,q\) of the form \(4n+1\). To establish the assumption,
we therefore only need to prove that the forms \((G)\) certainly
contain \emph{prime numbers}. This already seems very plausible
in itself, since these forms, together with \(4pqt+h,4pqt+h'\), and so on,
\SourceLeaf{362}
\((H)\), contain all numbers coprime to \(4pq\), and therefore
also all prime numbers (except \(2,p,q\)). There seems to be no reason
why the series of primes should not be equally distributed among
those forms, so that an eighth belong to \((G)\) and the rest to
\((H)\). Yet such reasoning clearly falls far short of geometric
rigor. The distinguished Le Gendre himself admits that it seems
quite difficult to prove the theorem that a form \(kt+l\), with
\(k,l\) given coprime numbers and \(t\) varying, certainly contains
primes. He briefly points to a method that might lead there. But
many preliminary investigations seem necessary to us before this
route, at least, could lead to a rigorous proof. — — As for the
other assumption (III, second method), that there is a prime \(r\)
of the form \(3n+3\) of which another given prime \(p\) of the form
\(4n+1\) is a nonresidue, the distinguished Le Gendre has added
nothing at all. We proved above (Article 129) that primes of which
\(p\) is a nonresidue certainly exist. But our method hardly seems
suited to showing the existence of such primes \textit{that are
also of the form} \(4n+3\), as is required here, though not in our
first proof. We can, however, readily prove the truth of this
assumption as follows. By Article 287 there is a positive genus
of binary forms of determinant \(-p\) with character \(3,4;\ Np\).
Let \((a,b,c)\) be such a form, with \(a\) odd, as we may assume.
Then \(a\) will be of the form \(4n+3\), and will either itself
be prime or at least have a prime factor \(r\) of the form \(4n+3\).
But \(-pRa\), and therefore also \(-pRr\), giving \(pNr\). Still,
we must carefully note that the propositions of Articles 263, 287
rely on the fundamental theorem. So it would be circular reasoning
to base any part of it on them. — — Finally, the assumption in
III of the first method is much more unsupported still, so there
is no need to add anything further about it here.

Let us add an observation on case V, which the preceding method
does not adequately prove, but which is readily settled by the
following argument. If we had \(pNq,qNp\) simultaneously there, then
\(-pRq,-qRp\), from which we readily derive that \(-1\) is a characteristic
number of the form \((p,0,q)\). So, by the theory of ternary forms,
this form can be represented by \(xx+yy+zz\). Let
\[ptt+quu=(\alpha t+\beta u)^2+(\alpha't+\beta'u)^2+(\alpha''t+\beta''u)^2\]
or
\[\begin{gathered}\alpha\alpha+\alpha'\alpha'+\alpha''\alpha''=p,\\\beta\beta+\beta'\beta'+\beta''\beta''=q,\\\alpha\beta+\alpha'\beta'+\alpha''\beta''=0\end{gathered}\]

\SourceLeaf{363}
Then equations 1 and 2 make all \(\alpha,\alpha',\alpha'',\beta,\beta',\beta''\) odd. But clearly the third
equation cannot then hold. — Case II can also be settled in much
the same way.
